% Cluster E, part 1: augmented atoms and the common-domain theorem
% Body-only LaTeX; uses only preamble macros and fixed cross-cluster labels.

\section{Augmented atoms and the common domain}\label{sec:packets}

The free-component criterion of Theorem~\ref{thm:free-component} is exact, and its
persistence clause makes failures permanent; but as a \emph{positive} test it is
expensive. Deciding whether one component of $\Omfree$ reaches every required atom
demands a global analysis of how each of the finitely many collars, obstacles, and
exclusions of the frozen inventory sits relative to all of the others, and
Example~\ref{ex:no-common-component} shows that no atom-by-atom inspection can
substitute for that analysis: individual accessibility of every atom is compatible
with an unrepairable separation. From a \emph{partial} package --- the situation in
which the criterion would actually be invoked, with the inventory delivered piecemeal
by an ambient proof program --- the exact free-domain test is therefore difficult to
verify directly \cite{MWman}.

The present section proves a sufficient criterion that is \emph{source-checkable}:
each of its clauses is verified against named, finite source data rather than against
the component structure of an open set. The idea is to reverse the direction of the
analysis. Instead of asking how the frozen set $Q$ decomposes the localization, we
package $Q$ itself into finitely many pairwise disjoint full continua --- each atom
augmented by the inventory pieces assigned to it, the remainder grouped into
separated-obstacle packets --- while reserving one free port for each original atom.
Once the completed packages are pairwise disjoint, sphere duality forces their common
complement to be \emph{connected}; connectedness, not component bookkeeping, then
delivers the common free component that Theorem~\ref{thm:free-component} demands.

The section has three parts. Definition~\ref{def:packet} fixes the packet structure
and its typed assignment record; Lemma~\ref{lem:packets} is the duality argument; and
Theorem~\ref{thm:common-domain} assembles the sufficient criterion. How the completed
packets are actually computed from raw assigned pieces, and how their defining clauses
are audited with persistent failure certificates, is the subject of
Section~\ref{sec:audit}; Example~\ref{ex:merger} there shows why the audit must be run
on the completions and not on the raw pieces.

\subsection{Filled hulls and augmented packets}\label{subsec:E-packet-def}

Recall from Section~\ref{sec:prelim} the filled hull of a compact plane set $K$:
\begin{equation}\label{eq:E-fill}
\Fill(K)\;=\;K\,\cup\,\bigcup\bigl\{\,B : B\ \text{a bounded component of}\
\C\setminus K\,\bigr\}
\end{equation}
(equation~(9.5) of \cite{MWman}). Since $K$ is compact, $\C\setminus K$ has exactly
one unbounded component $\Omega_\infty(K)$, and \eqref{eq:E-fill} says
\[
\Fill(K)\;=\;\C\setminus\Omega_\infty(K).
\]
Consequently $\Fill(K)$ is closed, and it is bounded because every closed round disk
containing $K$ has its exterior inside $\Omega_\infty(K)$; thus $\Fill(K)$ is compact.
Its complement is the single connected set $\Omega_\infty(K)$, so $\Fill(K)$ is
\emph{full}, and it is the smallest full compact set containing $K$: if $F\supset K$
is full and compact, then $\C\setminus F$ is a connected unbounded subset of
$\C\setminus K$, hence contained in $\Omega_\infty(K)$, which gives
$\Fill(K)\subset F$ after taking complements. Finally, if $K$ is a continuum, then so
is $\Fill(K)$: every bounded component $B$ of $\C\setminus K$ has
$\partial B\subset K$, so its closure $\overline B$ is a connected set meeting $K$,
and $\Fill(K)=K\cup\bigcup\overline B$ is a union of connected sets all meeting the
connected set $K$. The passage from $K$ to $\Fill(K)$ is what the manuscript calls
\emph{fullification}; its algebra --- and in particular the fact that interleaved
adjoin-and-fill operations compute a well-defined completion --- is developed in
Section~\ref{sec:audit}.

Throughout, the \emph{atoms} $A_1,\dots,A_s\Subset U$ are the pairwise disjoint full
continua of Section~\ref{sec:certificates} that the connector must join, and $Q$ is
the frozen set of Section~\ref{sec:criterion} (equation~(9.4) of \cite{MWman}): the
union of the atoms, the closed terminal protected collars, the separated obstacles,
the critical-value exclusion, and the localization-forbidden set, fixed once and for
all before any connector is chosen.

\begin{definition}[Augmented-atom packet; manuscript Definition 10.9]\label{def:packet}
For the frozen set $Q$, an \emph{augmented-atom packet} consists of full continua
\[
\widehat A_1,\dots,\widehat A_s,\
\widehat O_1,\dots,\widehat O_t\;\Subset\;U
\]
and pairwise disjoint port arcs $\alpha_i\colon[0,1]\to U$ such that:
\begin{enumerate}
\item every protected collar, obstacle, critical exclusion, and other component of
$Q$ has one typed assignment, either to a connected atom seed $A_i\cup E_i$ or to a
connected separated-obstacle seed $O_\ell$;
\item $\widehat A_i=\Fill(A_i\cup E_i)$ for every $i$, and
$\widehat O_\ell=\Fill(O_\ell)$ for every $\ell$;
\item the completed packets
$\widehat A_1,\dots,\widehat A_s,\widehat O_1,\dots,\widehat O_t$ are pairwise
disjoint and cover $Q$; and
\item $\alpha_i(0)=x_i\in A_i$, $\alpha_i\cap\widehat A_i=\{x_i\}$, and
\[
\alpha_i\bigl((0,1]\bigr)\;\subset\;
U\setminus\Bigl(\bigcup_{k}\widehat A_k\cup\bigcup_{\ell}\widehat O_\ell\Bigr),
\]
and $x_i$ is not the impression of a protected chain.
\end{enumerate}
The assignment records the component name, the proof-use stages, the attachment
trace, the return word, the selected sheet, and the absorb-or-separate decision.
\end{definition}

Some commentary on the four clauses. In clause~(1), $E_i$ denotes the union of the
inventory pieces assigned to the $i$-th atom; it may be empty, and the clause demands
that the resulting \emph{atom seed} $A_i\cup E_i$ be connected, so that an assigned
piece disjoint from its atom must be accompanied by connecting material assigned to
the same label. The assignment is \emph{typed}: it is a function from the finitely
many named constituents of the frozen inventory to the packet labels, and each
assigned piece carries the six recorded fields listed after clause~(4). Those fields
--- which component of the ambient package the piece is, at which proof-use stages it
is read, its attachment trace, its return word, its selected sheet, and whether the
decision taken for it was absorption or separation --- are bookkeeping data of the
ambient proof program; the planar topology of this section reads only the geometry of
the seeds. The record is nevertheless load-bearing: it is what makes a later change
of assignment an explicit, auditable re-decision rather than a silent repair, a point
Example~\ref{ex:merger} makes concrete.

In clause~(2), fullness of the completed packets is automatic: the seeds are compact
and connected, so their filled hulls are full continua by the discussion following
\eqref{eq:E-fill}. Clause~(3) is the substantive geometric hypothesis, and it is
stated for the \emph{completed} packets $\widehat A_i,\widehat O_\ell$, not for the
raw seeds: fullification can create intersections that the seeds do not have, and an
audit run at seed level certifies nothing (Example~\ref{ex:merger}). Coverage of $Q$
means that every item of the frozen inventory lies inside some completed packet, so
that outside the packets there is no frozen material at all --- this is the inclusion
that will place the common complement inside $\Omfree$. Clause~(4) is the port
clause. It is deliberately stronger than the port clause of
Definition~\ref{def:certificate}: the arc's positive-parameter part must avoid every
completed hull, whereas the arcs of Definition~\ref{def:certificate} avoid only the
obstacle set, the atoms, and the terminal collar tails, and the arc here must meet
even its own augmented atom $\widehat A_i$ only at the landing point $x_i$. The final requirement,
that $x_i$ not be the impression of a protected chain, keeps the port compatible with
the protected marking; it reappears in Section~\ref{sec:audit} as the requirement
that $x_i$ lie in the unprotected exposure locus, outside the forbidden impression
set $P_i$.

\subsection{Nonseparation of finite disjoint full families}\label{subsec:E-duality}

The engine of the common-domain theorem is a nonseparation statement for finitely
many disjoint full continua on the sphere. Here a compact set $K\subset S^2$ is
called full when $S^2\setminus K$ is connected; for a compact subset of the plane
this agrees with the planar notion, because adjoining the point at infinity to the
unbounded complementary component does not change connectivity, so a plane-full
compact set is sphere-full and conversely. The proof uses the compacta form of
Alexander duality \cite[pp.~256--257]{Hat02}:
\begin{equation}\label{eq:E-duality}
\widetilde H_0\bigl(S^2\setminus K;\Z\bigr)\;\cong\;\check H^1(K;\Z)
\end{equation}
for an arbitrary nonempty proper compact subset $K\subset S^2$ (equation~(9.6) of
\cite{MWman}). The \v{C}ech form, rather than the locally contractible special case,
is required here: the atoms and the assigned inventory pieces are arbitrary continua,
no local connectivity is assumed anywhere in this paper, and for such compacta
singular cohomology need not compute the group on the right of
\eqref{eq:E-duality}.

\begin{remark}\label{rem:E-cech}
The insistence on the \v{C}ech form is not pedantry but a structural feature of the
framework. Clause~(6) of Theorem~\ref{thm:realization} promises that no local
connectivity of any atom is ever required, and the packets of
Definition~\ref{def:packet} inherit whatever boundary pathology their atoms and
assigned pieces possess. \v{C}ech cohomology is the invariant that remains dual to
the complement's component count for every compact subset of the sphere, without any
tameness hypothesis; every duality argument below is stated so that it factors
through \eqref{eq:E-duality} alone.
\end{remark}

\begin{lemma}[Finite disjoint full packets; manuscript Lemma 10.10]\label{lem:packets}
If $K_0,\dots,K_N\subset S^2$ are pairwise disjoint nonempty full continua, then
\[
S^2\setminus\bigcup_{r=0}^{N}K_r
\]
is connected.
\end{lemma}

\begin{proof}
Write $X:=\bigcup_{r=0}^{N}K_r$. We may assume that $X$ is a proper subset of
$S^2$, so that \eqref{eq:E-duality} applies to $X$ and to each $K_r$: were $X=S^2$,
the sets $K_0,\dots,K_N$ would be finitely many pairwise disjoint nonempty closed
sets covering $S^2$; each would then be open as well, being the complement of the
union of the others, and connectedness of $S^2$ would force $N=0$ and $K_0=S^2$, in
which case the complement is empty and there is nothing to prove.

\emph{Step 1: each packet has vanishing first \v{C}ech cohomology.}
Fix $r$. Fullness of $K_r$ makes $S^2\setminus K_r$ connected, and it is nonempty
because $K_r$ is proper. Being open in $S^2$, the complement is locally path
connected, so connectedness makes it path connected and its reduced zeroth singular
homology vanishes. Hence \eqref{eq:E-duality}, applied to the nonempty proper
compact set $K_r$, gives
\[
\check H^1(K_r;\Z)\;\cong\;\widetilde H_0\bigl(S^2\setminus K_r;\Z\bigr)\;=\;0 .
\]

\emph{Step 2: disjoint neighborhoods split fine covers.}
The sets $K_0,\dots,K_N$ are finitely many pairwise disjoint compacta, so their
mutual distances in the spherical metric are positive; choose $\varepsilon>0$ smaller
than half the minimum of these distances and let $N_r$ be the open
$\varepsilon$-neighborhood of $K_r$. The $N_r$ are pairwise disjoint open sets with
$K_r\subset N_r$.

Call a cover of $X$ by relatively open subsets \emph{fine} when every member is
contained in one of the sets $N_r$. The fine covers are cofinal among all relatively
open covers of $X$: an arbitrary cover $\mathcal U$ is refined by the fine cover
$\{\,V\cap N_r : V\in\mathcal U,\ 0\le r\le N\,\}$, since
$X\subset\bigcup_rN_r$. Let $\mathcal V$ be a fine cover with nonempty members. Each
member lies in exactly one $N_r$, because the $N_r$ are pairwise disjoint; write
$\mathcal V_r$ for the members contained in $N_r$. Since
$X\cap N_r=K_r$, the collection $\mathcal V_r$ is a relatively open cover of $K_r$.
Two members taken from distinct collections $\mathcal V_r,\mathcal V_{r'}$ are
disjoint, lying in the disjoint sets $N_r,N_{r'}$; hence no simplex of the nerve of
$\mathcal V$ mixes the collections, and the nerve of $\mathcal V$ is the disjoint
union of the nerves of $\mathcal V_0,\dots,\mathcal V_N$. Simplicial cohomology of a
finite disjoint union is the direct sum, so
\[
H^1\bigl(\mathrm{nerve}\,\mathcal V;\Z\bigr)\;\cong\;
\bigoplus_{r=0}^{N}H^1\bigl(\mathrm{nerve}\,\mathcal V_r;\Z\bigr).
\]

\emph{Step 3: pass to the direct limit.}
\v{C}ech cohomology of $X$ is the direct limit of the nerve cohomologies over the
directed system of relatively open covers, and the fine covers are cofinal in that
system, so the limit may be computed over fine covers only. A refinement of a fine
cover is again fine after intersecting with the $N_r$, and its restriction to $N_r$
refines $\mathcal V_r$; the splitting of Step~2 is therefore compatible with the
refinement maps, and restriction to $N_r$ is cofinal among the covers of $K_r$
because any relatively open cover of $K_r$ extends to a fine cover of $X$ by
adjoining covers of the other packets. Since a direct limit commutes with the finite
direct sum, Step~1 gives
\[
\check H^1(X;\Z)\;\cong\;\bigoplus_{r=0}^{N}\check H^1(K_r;\Z)\;=\;0 .
\]

\emph{Step 4: apply duality a second time.}
The set $X$ is compact, nonempty, and proper, so \eqref{eq:E-duality} applies to $X$
itself:
\[
\widetilde H_0\bigl(S^2\setminus X;\Z\bigr)\;\cong\;\check H^1(X;\Z)\;=\;0 .
\]
The complement $S^2\setminus X$ is nonempty, and a nonempty space with vanishing
reduced zeroth homology is connected. This is the assertion.
\end{proof}

Lemma~\ref{lem:packets} is the disjoint counterpart of
Lemma~\ref{lem:tree-full}. There, unions of nonseparating continua meeting in
connected singletons were handled by the Janiszewski separation lemma, one attachment
at a time; here the continua are pairwise disjoint, the count is arbitrary, and the
duality argument dispenses with any hypothesis on how the packets face one another.
Nothing is assumed about the boundaries of the $K_r$ --- they may be nowhere locally
connected --- and no width, diameter, or separation estimate enters.

\subsection{The common-domain theorem}\label{subsec:E-common}

\begin{theorem}[Augmented-Atom Common-Domain Theorem; manuscript Theorem
10.11]\label{thm:common-domain}
Every augmented-atom packet determines one component of $\Omfree$ with simultaneous
singleton ports to every original atom. Hence it supplies a protected regulated
connector.
\end{theorem}

\begin{proof}
Let an augmented-atom packet be given, with completed packets
$\widehat A_1,\dots,\widehat A_s,\widehat O_1,\dots,\widehat O_t\Subset U$ and port
arcs $\alpha_1,\dots,\alpha_s$ as in Definition~\ref{def:packet}.

\emph{The exterior packet.}
Adjoin to the family the set
\[
K_0\;:=\;S^2\setminus U .
\]
The Jordan-disk hypothesis makes $K_0$ a full continuum: by the Jordan--Schoenflies
theorem the Jordan curve $\partial U$ divides the sphere into two open Jordan disks
with common boundary $\partial U$, the closure of the one containing $\infty$ is the
closed Jordan disk $S^2\setminus U$, hence a continuum, and its complement is the
connected set $U$, so $K_0$ is full in the sphere sense used in
Lemma~\ref{lem:packets}. Each completed packet is a plane-full continuum, hence
sphere-full as noted in Section~\ref{subsec:E-duality}, and compact containment
$\widehat A_i,\widehat O_\ell\Subset U$ makes every completed packet disjoint from
$K_0$; together with clause~(3) of Definition~\ref{def:packet}, the family
\[
K_0,\ \widehat A_1,\dots,\widehat A_s,\ \widehat O_1,\dots,\widehat O_t
\]
consists of $1+s+t$ pairwise disjoint nonempty full continua in $S^2$.

\emph{The augmented complement is connected.}
Lemma~\ref{lem:packets} applies to this family and makes its common complement
connected. Since $S^2\setminus K_0=U$, that complement is
\begin{equation}\label{eq:E-omega-aug}
\Omega_{\mathrm{aug}}\;:=\;
U\setminus\Bigl(\bigcup_{i}\widehat A_i\cup\bigcup_{\ell}\widehat O_\ell\Bigr)
\;=\;S^2\setminus\Bigl(K_0\cup\bigcup_i\widehat A_i\cup
\bigcup_\ell\widehat O_\ell\Bigr),
\end{equation}
a connected subset of $\C$, open because $U$ is open and the packets are compact. It
is nonempty: clause~(4) places the positive-parameter part of each port arc inside
it.

\emph{The augmented complement lies in one free component.}
Clause~(3) of Definition~\ref{def:packet} makes the completed packets cover $Q$, so
\[
\Omega_{\mathrm{aug}}
\;=\;U\setminus\Bigl(\bigcup_i\widehat A_i\cup\bigcup_\ell\widehat O_\ell\Bigr)
\;\subset\;U\setminus Q\;=\;\Omfree .
\]
A nonempty connected subset of the open set $\Omfree$ lies in exactly one of its
components; let $\Omega_0$ be the component containing $\Omega_{\mathrm{aug}}$. This
is the component that the packet determines.

\emph{The ports are simultaneous singleton ports of $\Omega_0$.}
By clause~(4), each port interior satisfies
\[
\alpha_i\bigl((0,1]\bigr)\;\subset\;
U\setminus\Bigl(\bigcup_k\widehat A_k\cup\bigcup_\ell\widehat O_\ell\Bigr)
\;=\;\Omega_{\mathrm{aug}}\;\subset\;\Omega_0 .
\]
Moreover $A_i\subset A_i\cup E_i\subset\widehat A_i$, so
\[
\alpha_i\cap A_i\;\subset\;\alpha_i\cap\widehat A_i\;=\;\{x_i\},
\]
and $x_i=\alpha_i(0)\in A_i$ gives $\alpha_i\cap A_i=\{x_i\}$: the arc meets its atom
exactly in the one free point $x_i$. The arcs are pairwise disjoint Jordan arcs by
Definition~\ref{def:packet}. Thus $\Omega_0$ has simultaneous singleton ports to
every original atom, in the sense of Section~\ref{sec:criterion}.

\emph{Conclusion.}
Theorem~\ref{thm:free-component} converts the simultaneous singleton ports of
$\Omega_0$ into a protected ribbon-routing certificate, and
Theorem~\ref{thm:realization} realizes that certificate by a compact connected full
carrier with singleton attachment traces, protected crosscut survival, and preserved
cyclic order. The packet therefore supplies a protected regulated connector.
\end{proof}

\begin{remark}\label{rem:E-common-domain}
Three remarks of the manuscript \cite{MWman} deserve emphasis. First, no local
connectivity or width estimate occurs anywhere in the argument: the packets may have
arbitrarily wild boundaries, the duality step of Lemma~\ref{lem:packets} is
insensitive to them, and no modulus, diameter, or incidence bound is invoked at any
point. Second, the construction is natural under an orientation-preserving conformal
homeomorphism of Jordan disks: the boundary extensions can be glued to a sphere
homeomorphism, and such a map carries filled packets to filled packets, ports to
ports, the common domain \eqref{eq:E-omega-aug} to the corresponding common domain,
and any selected connector to a connector for the transported data, exactly as in
Proposition~\ref{prop:naturality}; the commutation of such a sphere homeomorphism
with the operator $\Fill$ is proved in detail in the naturality paragraph of
Theorem~\ref{thm:audit}. Third, the spanning tree produced inside the common
component is not claimed canonical: Theorem~\ref{thm:free-component} supplies
\emph{some} embedded tree through the free component, and different choices give
different carriers with the same protected marking. Only the component $\Omega_0$,
the ports, and the survival data are determined by the packet.
\end{remark}

Theorem~\ref{thm:common-domain} consumes completed packets: it presupposes that the
filled hulls $\widehat A_i,\widehat O_\ell$ have already been computed from the typed
assignment and checked against clauses~(1)--(4). In practice the assignment delivers
raw pieces one at a time, and both the computation of the completions and the
verification of the clauses raise questions of their own --- whether the result
depends on the order of processing, and whether individually verified ports must be
untangled globally. The next section resolves both questions and equips the
verification with persistent failure certificates.
